\documentclass[11pt]{article}
\usepackage{amsmath,amssymb,amsthm}
\newtheorem{theorem}{Theorem}
\title{The exact recursive equal-moment invariant collapses exterior readouts}
\author{}\date{}
\begin{document}\maketitle
\paragraph{Status and scope.}
This is a complete written obstruction to one proposed recursive invariant,
not a lower bound for HeadComplexity. It is not Lean-checked. It assumes an
exact balanced binary tree of positive atomic measures, equal first-half
moments at every split, and one scalar bias per node. It does not cover
approximate moment matching, other profile corrections, unequal moment
budgets, signed residues, or arbitrary admissible heads.

Consider a full binary tree with $H=2^d$ leaves. Each leaf carries a
positive mass at one real pole, and all poles are distinct. A node measure
is the sum of its descendant leaf measures. At every node with $2M$ leaves,
assume its two child measures have equal moments of degrees $0,\ldots,M-1$.
All queries at that node lie strictly to the right of all its poles.

\begin{theorem}
Under these assumptions, the difference of the two child Stieltjes
transforms is either strictly positive and strictly decreasing throughout
the exterior query interval, or strictly negative and strictly increasing.
After adding any scalar bias, its strict signs on ordered queries have at
most one sign change. In particular an alternating three-query word is
impossible. Therefore this invariant cannot realize every word at a
four-query root, regardless of how small or target-dependent its pole gaps are.
\end{theorem}

\begin{proof}
The degree-zero equality at every split divides a node's mass equally
between its equally sized children. Hence every leaf has the same weight
$w>0$. At the node under consideration let $Q_L,Q_R$ be the monic
degree-$M$ polynomials of its two child pole sets. Equality of moments
of degrees $1,\ldots,M-1$ means equality of the corresponding power sums.
Newton's identities show that all coefficients of $Q_L,Q_R$ except
their constant terms agree. Thus
\[
 Q_L=R-b,\qquad Q_R=R+b
\]
for a monic degree-$M$ polynomial $R$ and a nonzero real $b$.
The constant is nonzero because the child pole sets are disjoint.
Exchange the children if necessary to arrange $b>0$.

The polynomial $R$ has $M$ real roots, all below the exterior interval.
For completeness, $R'=Q_L'=Q_R'$ has $M-1$ distinct real critical
points by Rolle's theorem. Since both $R-b$ and $R+b$ have $M$ distinct
real roots, every local maximum of $R$ is greater than $b$, and every
local minimum is less than $-b$. Consequently $R$ crosses zero once on
each of its $M$ monotonicity intervals. These roots lie between the
corresponding roots of $R-b$ and $R+b$, so are below the largest child pole.
The assertion for $M=1$ follows directly from its unique linear root.

Equal weights give the exact logarithmic-derivative identity
\[
 F(t)=S_L(t)-S_R(t)
 =w\left(\frac{Q_L'(t)}{Q_L(t)}-\frac{Q_R'(t)}{Q_R(t)}\right)
 =\frac{2wbR'(t)}{R(t)^2-b^2}.                         \tag{1}
\]
On the exterior interval, $R(t)>b$, $R'(t)>0$, and $R''(t)\geq0$.
Writing the roots of $R$ as $r_1,\ldots,r_M$ gives
\[
 RR''-(R')^2=-R^2\sum_i(t-r_i)^{-2}<0.
\]
Therefore
\[
 R''(R^2-b^2)-2R(R')^2
 =R\bigl(RR''-2(R')^2\bigr)-b^2R''<0.
\]
Differentiating (1) proves $F>0$ and $F'<0$. Exchanging the children
negates $F$, proving the other orientation. A strictly monotone function
plus a constant has at most one sign change on ordered sample points.
\end{proof}

\paragraph{The remaining parameter count.}
If the first $H$ root moments, of degrees $0,\ldots,H-1$, are also
fixed, the common leaf weight is fixed and Newton's identities determine
the degree-$H$ pole polynomial up to its constant term. They do not fix
that final constant. Thus there is at most one continuous pole parameter,
not the $H$-dimensional positive moment fiber available before imposing
all descendant equalities.

For a fixed pole multiset, the hierarchical partition is unique up to
exchanging the two children at each internal node. Indeed the monic
polynomial square part $R$ in $Q=R^2-b^2$ is unique: if two monic
degree-$M$ polynomials had squares differing by a constant, their
difference times their sum could only be constant if the difference
were zero. Repeating down the tree yields at most $2^{H-1}$ ordered
labelings. The same argument shows that the hierarchy is generated by
iterated quadratic polynomial preimages. These observations explain
why independent local moment-fiber dimensions cannot be added over levels.

\paragraph{Relation to the local split theorem.}
The theorem in \texttt{positive-moment-split.tex} remains valid: its two
children may have independently varying positive weights. Imposing all
lower-level equal-mass conditions forces those leaf weights to agree,
and then (1) removes the exterior sign freedom. A viable use of positive
moment geometry must relax or replace at least one part of this exact
recursive invariant; the local inverse-function theorem alone cannot
provide the missing compatibility.
\end{document}
